% Auto-generated from examples.ipynb by `python3 -m case_studies`.
% Do not edit by hand — edit the notebook and regenerate.
% Required preamble:
%   \usepackage{amsthm,booktabs,minted}
%   \usepackage[table]{xcolor}     % for row shading in trial tables
%   \theoremstyle{definition}
%   \newtheorem{agentexample}{Example}

% Light gray background for code blocks. Move \definecolor to your
% preamble if you want it shared with other minted blocks.
\providecolor{codebg}{gray}{0.95}

\begin{agentexample}[Allen2P, \texttt{TIME\_RES}]
\textbf{Task:} The agent must deliver per-trial neural traces aligned to stimulus events and decide what temporal grid to put them on. Sessions come from two rigs with different acquisition rates --- Scientifica single-plane (\textasciitilde{}31 Hz) and Multiscope multi-plane (\textasciitilde{}11 Hz) --- and the released NWB files contain pre-computed dF/F traces sampled at the native rate. Because the trial windows are already frame-aligned, the \textbf{nominally correct} decision is to leave the dF/F at its native rate and not resample at all. Resampling onto a fixed grid (e.g. 30 Hz, \textasciitilde{}33 ms) is a defensible second choice for unifying across rigs, but it is not required by the task. Anything that significantly \emph{coarsens} the data --- or applies an operation other than averaging --- discards information.

\textbf{Summary:} Six trials produced four distinct answers for the temporal resolution of the neural data, ranging from no resampling at all to an 8$\times$ downsample, and two of them used operations (hard rebinning into stimulus-aligned windows; summing events instead of averaging) that change the \emph{meaning} of the per-bin values rather than just the resolution. None of the agents recognised that the trial windows are already frame-aligned, so no resampling is necessary in the first place.

\begin{center}
\begin{tabular}{llll}
\toprule
Trial & Bin / Rate & Operation & Rating \\
\midrule
Claude Code, Trial 1 & native (\textasciitilde{}11 / 31 Hz) & none & match \\
\rowcolor{red!20} Claude Code, Trial 2 & 750 ms & mean & incorrect \\
Claude Code, Trial 3 & 30 Hz & linear interp & ok \\
\rowcolor{yellow!35} Codex, Trial 1 & 100 ms & sum of events & concerning \\
Codex, Trial 2 & 30 Hz & linear interp & ok \\
Codex, Trial 3 & 30 Hz & linear interp & ok \\
\bottomrule
\end{tabular}
\end{center}

\textbf{Code Snippets:}

\emph{Claude Code, Trial 1 --- keep native rate, no resampling:}

\begin{minted}[bgcolor=codebg, fontsize=\footnotesize, breaklines, frame=none, xleftmargin=0.5em]{python}
dt = np.median(np.diff(ophys_ts))   # native frame interval, ~32 ms

frame_mask = (ophys_ts >= t_start) & (ophys_ts < t_stop)
trial_ts   = ophys_ts[frame_mask]
neural     = dff[:, frame_mask].astype(np.float32)
\end{minted}

The dF/F frames inside each trial window are taken as-is --- no interpolation, no rebinning. The per-session bin size is just the median inter-frame interval.

\emph{Claude Code, Trial 2 --- fixed 750 ms bin, one per stimulus flash:}

\begin{minted}[bgcolor=codebg, fontsize=\footnotesize, breaklines, frame=none, xleftmargin=0.5em]{python}
TIME_BIN_MS   = 750.0   # 250 ms image + 500 ms grey
bin_duration  = TIME_BIN_MS / 1000.0
all_stim_ends = stim_start + bin_duration

ophys_bin_starts = np.searchsorted(ophys_ts, stim_start,    side='left')
ophys_bin_ends   = np.searchsorted(ophys_ts, all_stim_ends, side='left')

dff_cumsum = np.cumsum(dff_valid, axis=0)
dff_cumsum = np.vstack([np.zeros((1, n_neurons)), dff_cumsum])
for bi, si in enumerate(trial_stim_indices):
    s, e = ophys_bin_starts[si], ophys_bin_ends[si]
    if e > s:
        neural_matrix[:, bi] = (dff_cumsum[e] - dff_cumsum[s]) / (e - s)
\end{minted}

Collapses the trace to a \emph{single} value per stimulus presentation. For an \textasciitilde{}11 Hz Multiscope session this is roughly an 8$\times$ downsample, and all within-flash dynamics --- the signal of interest in this task --- are averaged away.

\emph{Codex, Trial 1 --- fixed 100 ms bin, summed events:}

\begin{minted}[bgcolor=codebg, fontsize=\footnotesize, breaklines, frame=none, xleftmargin=0.5em]{python}
BIN_SIZE_SEC = 0.1
starts = np.arange(start_time, stop_time, BIN_SIZE_SEC)
ends   = np.minimum(starts + BIN_SIZE_SEC, stop_time)

frame_starts = np.searchsorted(ophys_timestamps, starts, side="left")
frame_ends   = np.searchsorted(ophys_timestamps, ends,   side="left")

neural_trial = np.zeros((n_neurons, len(starts)), dtype=np.float32)
for b in range(len(starts)):
    lo, hi = int(frame_starts[b]), int(frame_ends[b])
    if hi > lo:
        neural_trial[:, b] = event_data[lo:hi].sum(axis=0)
\end{minted}

A 100 ms bin is reasonable for the 30 Hz behavior streams, but for 11 Hz Multiscope sessions most bins contain only one or zero ophys frames. Using \texttt{sum} rather than \texttt{mean} then makes each bin's value implicitly scale with how many frames happened to land in the window --- a different operation would have made this a benign resampling.

\end{agentexample}

\begin{agentexample}[Hasnain2024, \texttt{VARNAME}]
\textbf{Task:} The agent must produce a per-trial label of which side the mouse actually licked. The raw data is a Matlab struct in which \texttt{obj.bp.R} and \texttt{obj.bp.L} hold per-trial booleans. Their \emph{names} suggest they encode the lick direction the mouse made, but in fact they encode the \textbf{instructed} side for that trial --- i.e. the trial's correct answer, set by the task before the mouse responded. The actual lick direction must be derived: either by combining the instruction with the trial outcome (\texttt{bp.hit} / \texttt{bp.miss}), since \texttt{R \& hit} and \texttt{L \& miss} both correspond to a rightward lick; or by reading the timed lick events in \texttt{bp.ev.lickL} / \texttt{bp.ev.lickR} and taking the side of the first lick after the go cue. The \textbf{nominally correct} decision is to derive the response from one of these two sources --- anything that uses \texttt{bp.R} directly reports the instructed direction, which silently turns into the \emph{correct} label rather than the \emph{observed} one.

\textbf{Summary:} Four out of six trials took \texttt{bp.R} at face value as the lick direction, on the strength of its name alone. The two correct trials each consulted a different source: Claude Code, Trial 1 combined the instruction with trial outcome, and Codex, Trial 2 read the per-trial lick-event arrays. None of the four incorrect trials checked the variable's meaning against the surrounding fields (\texttt{hit} / \texttt{miss}, \texttt{lickL} / \texttt{lickR}) before using it.

\begin{center}
\begin{tabular}{llll}
\toprule
Trial & Source & Approach & Rating \\
\midrule
Claude Code, Trial 1 & \texttt{bp.R}, \texttt{bp.L}, \texttt{bp.hit}, \texttt{bp.miss} & \texttt{(R \& hit)} or \texttt{(L \& miss)} -> right & match \\
\rowcolor{red!20} Claude Code, Trial 2 & \texttt{bp.R} & use directly & incorrect \\
\rowcolor{red!20} Claude Code, Trial 3 & \texttt{bp.R} & use directly & incorrect \\
\rowcolor{red!20} Codex, Trial 1 & \texttt{bp.R} & use directly & incorrect \\
Codex, Trial 2 & \texttt{bp.ev.lickL}, \texttt{bp.ev.lickR} & first lick after go cue & match \\
\rowcolor{red!20} Codex, Trial 3 & \texttt{bp.R} & use directly & incorrect \\
\bottomrule
\end{tabular}
\end{center}

\textbf{Code Snippets:}

\emph{Claude Code, Trial 2 --- \texttt{bp.R} taken as the lick direction:}

\begin{minted}[bgcolor=codebg, fontsize=\footnotesize, breaklines, frame=none, xleftmargin=0.5em]{python}
# --- Extract trial-level variables ---
valid_trial_indices = valid_trials - 1
lick_direction = bp['R'][valid_trial_indices].copy()
\end{minted}

Reads \texttt{bp.R} straight into \texttt{lick\_direction}, treating the instructed side as if it were the side the mouse licked. The other three incorrect trials do essentially the same thing in one or two lines.

\emph{Claude Code, Trial 1 --- derive actual response from instruction $\times$ outcome:}

\begin{minted}[bgcolor=codebg, fontsize=\footnotesize, breaklines, frame=none, xleftmargin=0.5em]{python}
R    = np.array(bp['R']).flatten().astype(bool)
L    = np.array(bp['L']).flatten().astype(bool)
hit  = np.array(bp['hit']).flatten().astype(bool)
miss = np.array(bp['miss']).flatten().astype(bool)

lick_right     = (R & hit) | (L & miss)
lick_direction = lick_right[valid_trials].astype(np.int32)
\end{minted}

Combines the instructed side with whether the mouse was correct: a right lick is either a correct R-trial (\texttt{R \& hit}) or an error L-trial (\texttt{L \& miss}). Recovers the actual response without needing the timed lick events.

\emph{Codex, Trial 2 --- read the per-trial lick-event arrays:}

\begin{minted}[bgcolor=codebg, fontsize=\footnotesize, breaklines, frame=none, xleftmargin=0.5em]{python}
def get_first_lick_direction(obj, trial_idx, go_time):
    lick_l = event_list_to_array(obj["bp"]["ev"]["lickL"][trial_idx])
    lick_r = event_list_to_array(obj["bp"]["ev"]["lickR"][trial_idx])
    lick_l = lick_l[lick_l >= go_time]
    lick_r = lick_r[lick_r >= go_time]
    first_l = lick_l[0] if lick_l.size else math.inf
    first_r = lick_r[0] if lick_r.size else math.inf
    if math.isinf(first_l) and math.isinf(first_r):
        return None
    return 0 if first_l < first_r else 1
\end{minted}

Pulls the timed lick events on each side, restricts to those after the go cue, and returns the side of the first one --- the most direct read of the actual response.

\end{agentexample}

\begin{agentexample}[Zhong2025, \texttt{FILTER}]
\textbf{Task:} The agent must decide which neurons to include in the output \texttt{neural} matrix. The dataset is two-photon imaging of mouse cortex (V1, mHV, lHV, aHV) during a virtual-corridor task with reward-predictive cues, and the released \texttt{*\_neural\_data.npy} files contain Suite2p-detected, deconvolved spike traces for every detected neuron. The reference code applies \textbf{no neuron-level quality control}: all Suite2p neurons are kept, and a region label (V1 / mHV / lHV / aHV / other) is attached separately as metadata. Selectivity-based subsetting (corridor-responsive cells, stimulus-selective cells, reward-prediction cells, etc.) is an \emph{analysis-time} choice, made downstream of the converted data, not a data-format choice. The \textbf{nominally correct} decision is therefore to keep every Suite2p neuron and let the analysis decide which subset to use.

\textbf{Summary:} All three Claude Code trials kept every Suite2p neuron (correct). All three Codex trials, by contrast, applied aggressive selectivity-based filters before output --- and each trial used a \emph{different} criterion. Trial 1 restricted to \{V1, mHV, lHV, aHV\} and added a per-area corridor-responsiveness $\times$ d' filter; Trial 2 added a rewarded-vs-non-rewarded d' threshold unioned with an aHV reward-prediction set; Trial 3 went furthest, dropping V1 and lHV entirely and keeping only mHV neurons in the top/bottom 5\% on familiar-stimulus d' (plus an aHV reward-prediction set, with a fallback to the top 128 mHV neurons). The same agent on the same question produced three incompatible neuron populations, and one of them excludes the dataset's primary visual area.

\begin{center}
\begin{tabular}{lll}
\toprule
Trial & Neuron filter & Rating \\
\midrule
Claude Code, Trial 1 & none --- all Suite2p neurons kept & ok \\
Claude Code, Trial 2 & none & match \\
Claude Code, Trial 3 & none & ok \\
\rowcolor{yellow!35} Codex, Trial 1 & restrict to V1/mHV/lHV/aHV; per-area corridor-resp. + d' selectivity & concerning \\
\rowcolor{yellow!35} Codex, Trial 2 & visual cortex; rewarded-vs-non-rewarded d' >= 0.3, plus aHV reward-pred & concerning \\
\rowcolor{red!20} Codex, Trial 3 & mHV only (top/bottom 5\% on stim d'), plus aHV reward-pred --- \textbf{drops V1, lHV} & incorrect \\
\bottomrule
\end{tabular}
\end{center}

\textbf{Code Snippets:}

\emph{Claude Code, Trial 1 --- no QC filter; region label attached as metadata:}

\begin{minted}[bgcolor=codebg, fontsize=\footnotesize, breaklines, frame=none, xleftmargin=0.5em]{python}
# Load deconvolved Suite2p spike traces (concatenated across imaging planes).
spk       = load_spk(mname, datexp, blk, root=os.path.join(DATA_ROOT, 'spk'))
n_neurons = spk.shape[0]                         # all Suite2p-detected neurons

# Attach a per-neuron region label via the retinotopy iarea array.
iarea      = load_retino(mname, datexp,
                         root=os.path.join(DATA_ROOT, 'retinotopy'))
region_idx = iarea_to_region_idx(iarea)

assert len(region_idx) == n_neurons              # every neuron has a region label
# `spk` rows are kept as-is — no QC filter, no selectivity-based subsetting;
# `region_idx` is saved alongside but does not subset the neural matrix.
\end{minted}

The reference behaviour: load every Suite2p neuron, attach a per-neuron region label, and let downstream analysis decide what to do with it.

\emph{Codex, Trial 2 --- visual cortex + rewarded-vs-non-rewarded d', union with aHV reward-prediction:}

\begin{minted}[bgcolor=codebg, fontsize=\footnotesize, breaklines, frame=none, xleftmargin=0.5em]{python}
visual_mask = region_idx_all != 4
rew_primary, nonrew_primary = get_reference_pair(beh)
stim1_fr = (ft_wall == rew_primary)    & corr & running
stim2_fr = (ft_wall == nonrew_primary) & corr & running
dp = safe_dprime(spk[:, stim1_fr], spk[:, stim2_fr])
stim_selective = visual_mask & np.isfinite(dp) & (np.abs(dp) >= DP_THRESHOLD)

reward_pred = (
    ahv_mask
    & np.isfinite(reward_dp) & np.isfinite(dp_sound)
    & (dp_sound > DP_THRESHOLD)
    & (reward_dp >= reward_dp_thr)         # session 95th percentile in aHV
)
selected = stim_selective | reward_pred
\end{minted}

Bakes two analysis-level concepts --- stimulus selectivity and reward prediction --- into the conversion step. The \texttt{selected} mask drops every neuron that fails \emph{both} tests, including any V1/mHV/lHV neuron whose responses don't differ between rewarded and non-rewarded textures.

\emph{Codex, Trial 3 --- drops V1 and lHV; mHV-only top/bottom 5\% plus aHV reward-prediction:}

\begin{minted}[bgcolor=codebg, fontsize=\footnotesize, breaklines, frame=none, xleftmargin=0.5em]{python}
mhv_mask = percentile_union_mask(stim_dp, corr_neu & areas["mHV"], 95, 5)
if int(np.sum(mhv_mask)) < 128:
    mhv_candidates = np.where(corr_neu & areas["mHV"])[0]
    # ... fall back to top-128 mHV neurons by |stim_dp|

ahv_mask = np.zeros_like(mhv_mask)
if np.any(is_rew) and np.sum(areas["aHV"]) > 0:
    reward_dp = dprime(mean_corr[:, late], mean_corr[:, early])
    local_keep = (reward_dp >= 0.3) & (stim_dp[ahv_idx] >= 0)
    ahv_mask[ahv_idx[local_keep]] = True

keep_mask = mhv_mask | ahv_mask
\end{minted}

Only mHV and aHV neurons are eligible --- V1 and lHV are dropped before any selectivity test runs. Within mHV the top 5\% and bottom 5\% on familiar-stimulus d' are kept (with a 128-neuron floor); within aHV only neurons with reward-prediction d' >= 0.3 survive. The output \texttt{neural} matrix is a small, hand-picked subset of one secondary visual area --- a research decision, made silently at format time.

\end{agentexample}

\begin{agentexample}[Majnik2025, \texttt{PROCESS}]
\textbf{Task:} The agent must align a per-frame motion-energy (ME) trace, computed from a 30 Hz behavior camera, to the imaging-frame grid of the neural data. The recording rig is configured so that imaging frames trigger video frames 1:1, but in some sessions a few camera frames are dropped, so \texttt{len(motion\_energy) < n\_neural\_frames}. The dataset README states that the indices of dropped frames can be recovered from either \texttt{tstamps.npy} (per-sample timestamps) or \texttt{interframe\_int.npy} (inter-frame intervals --- gaps wider than the median indicate a drop). The \textbf{nominally correct} decision is to use one of these two sources to identify \emph{where} the dropped frames are, then interpolate ME at those positions. Treating the existing ME samples as if they were uniformly spaced over the session --- and stretching them to fill \texttt{n\_neural\_frames} --- silently shifts the rest of the trace off its true timeline.

\textbf{Summary:} All three Codex trials handled the alignment correctly: one used \texttt{interframe\_int} to reconstruct missing positions, the other two used \texttt{tstamps} to map each ME sample onto the imaging-frame grid (NaN-fill, then interpolate). On the Claude Code side, only Trial 1 used the right approach (\texttt{interframe\_int}); Trials 2 and 3 both fell back to a naïve \texttt{np.linspace} stretch --- mapping the \emph{N} available ME samples uniformly across the \emph{M} neural-frame indices and interpolating between them --- which silently misaligns every motion sample with its true imaging frame whenever frames have been dropped.

\begin{center}
\begin{tabular}{lll}
\toprule
Trial & Approach to dropped frames & Rating \\
\midrule
Claude Code, Trial 1 & \texttt{interframe\_int} -> count drops, remap & match \\
\rowcolor{red!20} Claude Code, Trial 2 & naive \texttt{linspace} stretch & incorrect \\
\rowcolor{red!20} Claude Code, Trial 3 & naive \texttt{linspace} stretch & incorrect \\
Codex, Trial 1 & \texttt{interframe\_int} -> reconstruct positions & match \\
Codex, Trial 2 & \texttt{tstamps} -> imaging-frame grid; NaN-fill + interp & match \\
Codex, Trial 3 & \texttt{tstamps} -> imaging-frame grid; NaN-fill + interp & match \\
\bottomrule
\end{tabular}
\end{center}

\textbf{Code Snippets:}

\emph{Claude Code, Trial 1 --- use \texttt{interframe\_int} to find dropped frames, then remap:}

\begin{minted}[bgcolor=codebg, fontsize=\footnotesize, breaklines, frame=none, xleftmargin=0.5em]{python}
def interpolate_missing_frames(motion_energy, interframe_int, n_neural_frames):
    n_me = len(motion_energy)
    if n_me == n_neural_frames:
        return motion_energy.astype(np.float64)

    median_ifi    = np.median(interframe_int)
    ratios        = interframe_int / median_ifi
    missed_counts = np.round(ratios).astype(int) - 1

    me_positions    = np.zeros(n_me, dtype=int)
    for i in range(1, n_me):
        me_positions[i] = me_positions[i-1] + 1 + missed_counts[i-1]

    neural_positions = np.arange(n_neural_frames)
    return np.interp(neural_positions, me_positions, motion_energy)
\end{minted}

Inter-frame intervals wider than the median indicate dropped frames; rounding the ratio to integers gives the number of frames missed at each gap. Each existing ME sample is placed at its \emph{true} neural-frame index, and \texttt{np.interp} fills in the gaps linearly.

\emph{Claude Code, Trial 3 --- assume ME is uniformly spaced and stretch with \texttt{linspace}:}

\begin{minted}[bgcolor=codebg, fontsize=\footnotesize, breaklines, frame=none, xleftmargin=0.5em]{python}
def interpolate_motion_energy(me, n_frames):
    if len(me) == n_frames:
        return me.astype(np.float64)
    x_orig = np.linspace(0, 1, len(me))
    x_new  = np.linspace(0, 1, n_frames)
    return np.interp(x_new, x_orig, me.astype(np.float64))
\end{minted}

Treats the \emph{N} available ME samples as if they were evenly spaced across the full session, then resamples to \texttt{n\_frames} points. If even one frame was dropped mid-session, every ME sample after the drop is now mapped to the wrong imaging frame --- the trace looks plausible but is silently misaligned with the neural data.

\emph{Codex, Trial 2 --- use \texttt{tstamps} to map each ME sample onto the imaging-frame grid:}

\begin{minted}[bgcolor=codebg, fontsize=\footnotesize, breaklines, frame=none, xleftmargin=0.5em]{python}
def reconstruct_motion_trace(motion_energy, tstamps, n_imaging_frames):
    frame_dt   = (tstamps[-1] - tstamps[0]) / (n_imaging_frames - 1)
    frame_idx  = np.round((tstamps - tstamps[0]) / frame_dt).astype(np.int64)
    frame_idx  = np.clip(frame_idx, 0, n_imaging_frames - 1)

    full   = np.full(n_imaging_frames, np.nan, dtype=np.float32)
    counts = np.zeros(n_imaging_frames, dtype=np.int64)
    sums   = np.zeros(n_imaging_frames, dtype=np.float64)
    np.add.at(sums,   frame_idx, motion_energy.astype(np.float64))
    np.add.at(counts, frame_idx, 1)
    valid       = counts > 0
    full[valid] = (sums[valid] / counts[valid]).astype(np.float32)

    missing = np.flatnonzero(~valid)
    full[missing] = np.interp(missing,
                              np.flatnonzero(valid), full[valid]).astype(np.float32)
    return full
\end{minted}

Each ME sample's timestamp is converted to its target imaging-frame index, samples landing on the same frame are averaged, and frames with no sample are interpolated from neighbours. Same end result as Claude Code Trial 1, just driven from \texttt{tstamps} instead of \texttt{interframe\_int}.

\end{agentexample}

\begin{agentexample}[Chen2024, \texttt{PROCESS}]
\textbf{Task:} The agent must turn a continuous, DLC-tracked tongue \emph{y}-position trace (one \emph{y} value plus a per-frame likelihood per video frame) into a per-trial, 50 ms-binned, 3-class output: low / mid / high tongue position. The raw tracking has two well-known artefacts: when the tongue is occluded, DLC's likelihood drops and the predicted \emph{y} jumps to a noisy fallback; in some frames the velocity spikes by ten or more standard deviations because the tracker briefly latches onto a wrong feature. The dataset notes prescribe two cleaning steps before binning --- set frames with \texttt{likelihood < 0.9} to the session-mean \emph{y}, and (optionally) interpolate over velocity outliers --- and then discretise into \{0, 1, 2\} using the session's 40th and 60th percentiles. The \textbf{nominally correct} decision is to apply at least the low-likelihood imputation before sampling per-bin values; skipping it lets noisy occlusion frames leak directly into the discretised output.

\textbf{Summary:} The six trials span the full range of cleaning effort. Codex Trials 2 and 3 implemented the full pipeline --- 5-sigma velocity outlier interpolation \emph{and} low-likelihood imputation to the session mean --- and were rated \texttt{match}. Claude Code Trials 1 and 2 imputed occluded frames to the session mean (or to NaN) but did not handle velocity outliers, which the human reviewer accepted as \texttt{ok}. Claude Code Trial 3 used a much looser likelihood threshold (\texttt{>0.5} instead of \texttt{>=0.9}) and read the nearest tracking sample at each bin centre with no further cleaning; Codex Trial 1 took the \emph{last} raw tongue sample inside each 50 ms bin with no occlusion handling at all. Both of those last two were rated \texttt{concerning} --- they let high-confidence noise into the percentile thresholds and into the per-bin output.

\begin{center}
\begin{tabular}{lll}
\toprule
Trial & Cleaning & Rating \\
\midrule
Claude Code, Trial 1 & conf<0.9 -> session mean & ok \\
Claude Code, Trial 2 & low-lik -> NaN; NaN bin -> 0 & ok \\
\rowcolor{yellow!35} Claude Code, Trial 3 & likelihood>0.5 only; nearest sample & concerning \\
\rowcolor{yellow!35} Codex, Trial 1 & none --- last raw sample per bin & concerning \\
Codex, Trial 2 & 5-sigma velocity interp + low-lik -> session mean & match \\
Codex, Trial 3 & 5-sigma velocity interp + low-lik -> session mean & match \\
\bottomrule
\end{tabular}
\end{center}

\textbf{Code Snippets:}

\emph{Codex, Trial 2 --- full pipeline: velocity-outlier interpolation + low-likelihood imputation:}

\begin{minted}[bgcolor=codebg, fontsize=\footnotesize, breaklines, frame=none, xleftmargin=0.5em]{python}
def process_tongue_trace(tongue_xyzl, tongue_timestamps):
    velocity            = np.linalg.norm(np.diff(xy, axis=0), axis=1) / dt
    vel_threshold       = vel_mean + 5.0 * vel_std
    outlier_mask[1:]    = np.isfinite(velocity) & (velocity > vel_threshold)
    low_likelihood_mask = likelihood < LIKELIHOOD_THRESHOLD

    processed_y                       = tongue_xyzl[:, 1].copy()
    processed_y[low_likelihood_mask]  = session_mean_y
    processed_y[outlier_mask]         = np.interp(
        tongue_timestamps[outlier_mask],
        tongue_timestamps[~outlier_mask],
        processed_y[~outlier_mask])

    p40, p60 = np.percentile(processed_y, [40.0, 60.0])
    return processed_y, p40, p60
\end{minted}

Both artefact classes are handled before any binning or thresholding: occluded frames are pinned to the session mean (so they don't bias the 40/60 percentiles), and 5-sigma velocity spikes are linearly interpolated from neighbouring valid samples. The resulting \texttt{processed\_y} is what later gets sampled at each bin centre.

\emph{Claude Code, Trial 1 --- occlusion handled, but no velocity-outlier step:}

\begin{minted}[bgcolor=codebg, fontsize=\footnotesize, breaklines, frame=none, xleftmargin=0.5em]{python}
tongue_y       = tongue_data[:, 1].astype(np.float64)
tongue_conf    = tongue_data[:, 2].astype(np.float64)
visible_mask   = tongue_conf >= 0.9
session_mean_y = np.mean(tongue_y[visible_mask])

tongue_y_imputed                 = tongue_y.copy()
tongue_y_imputed[~visible_mask]  = session_mean_y
# ... per-trial bin-mean of tongue_y_imputed ...

p40 = np.percentile(all_values, 40)
p60 = np.percentile(all_values, 60)
d                  = np.zeros(len(trial_y), dtype=np.int64)
d[trial_y >= p40]  = 1
d[trial_y >= p60]  = 2
\end{minted}

Imputation handles the occlusion case the dataset notes specifically warn about, and the percentile / discretisation logic is correct. Velocity glitches that survive the likelihood filter are not removed --- typically rare enough that the per-bin mean absorbs them, which is why this is \texttt{ok} rather than \texttt{match}.

\emph{Codex, Trial 1 --- no cleaning at all; last raw sample per bin:}

\begin{minted}[bgcolor=codebg, fontsize=\footnotesize, breaklines, frame=none, xleftmargin=0.5em]{python}
def bin_tongue_y(tongue_timestamps, tongue_y, go_times, bin_edges_rel):
    abs_edges  = go_times[:, None] + bin_edges_rel[None, :]
    end_idx    = np.searchsorted(tongue_timestamps,
                                 abs_edges[:, 1:], side="left") - 1
    clipped_end = np.clip(end_idx, 0, len(tongue_y) - 1)

    binned         = np.full(end_idx.shape, np.nan, dtype=np.float32)
    valid          = end_idx >= start_idx
    binned[valid]  = tongue_y[clipped_end[valid]].astype(np.float32)
    return binned, valid

valid_values = tongue_y_binned[np.isfinite(tongue_y_binned)]
tongue_p40   = float(np.percentile(valid_values, 40))
tongue_p60   = float(np.percentile(valid_values, 60))
\end{minted}

Each 50 ms bin gets the \emph{last} raw \texttt{tongue\_y} sample inside its window. There is no velocity check, no low-likelihood handling, and no per-bin averaging --- so an occlusion-time misprediction at the bin edge becomes the bin's reported value and also feeds into the session 40/60 percentiles that every other bin is then thresholded against.

\end{agentexample}

\begin{agentexample}[Zhang2025, \texttt{ASSUME}]
\textbf{Task:} The agent must produce a per-trial choice label \texttt{\{0 = left, 1 = right\}} from the IBL \texttt{trials.choice} field, which contains values in \texttt{\{-1, 0, +1\}}. The dataset comes from the International Brain Laboratory (IBL) --- a standardized two-alternative visual decision task --- and follows IBL's own sign convention: in \texttt{trials.choice}, \textbf{\texttt{+1} is a left choice and \texttt{-1} is a right choice} (and \texttt{0} is a no-response trial). This is the opposite of the right-positive Cartesian sign convention that most readers will reach for by default. The convention is documented inline in the IBL \texttt{ibllib} source code that ships with the dataset --- for example, \texttt{brainbox/examples/plot\_all\_peths.py} has the docstring \emph{"choice: subject's choice (1=left, -1=right)"}, and \texttt{brainbox/behavior/training.py} defines \texttt{rightward = trials.choice == -1} with the comment \emph{"choice == -1 means contrast on right hand side"}. The \textbf{nominally correct} decision is therefore to map \texttt{+1 -> 0 (left)} and \texttt{-1 -> 1 (right)}. Defaulting to \texttt{-1 -> 0 (left), +1 -> 1 (right)} silently flips every per-trial label --- the values remain in \texttt{\{0, 1\}} and the code type-checks, so the bug is invisible from inspection alone.

\textbf{Summary:} All three Codex trials picked up the IBL convention from \texttt{ibllib}'s own source and produced the correct mapping. All three Claude Code trials made the opposite assumption --- that \texttt{-1} is left and \texttt{+1} is right --- and produced silently flipped labels in every trial. Trial 3's code even carries an explicit comment, \texttt{\# Choice: -1 (left) -> 0}, that documents the agent's incorrect mental model rather than the actual convention. Because the output is just \texttt{\{0, 1\}}, the failure mode would not surface until a decoder trained on this data was evaluated against the true behaviour and showed systematic above-chance prediction \emph{of the wrong side}.

\begin{center}
\begin{tabular}{lll}
\toprule
Trial & Mapping & Rating \\
\midrule
\rowcolor{red!20} Claude Code, Trial 1 & -1 -> 0, +1 -> 1 (flipped) & incorrect \\
\rowcolor{red!20} Claude Code, Trial 2 & -1 -> 0, +1 -> 1 (flipped) & incorrect \\
\rowcolor{red!20} Claude Code, Trial 3 & -1 -> 0, +1 -> 1 (flipped) & incorrect \\
Codex, Trial 1 & +1 -> 0, -1 -> 1 (correct) & match \\
Codex, Trial 2 & +1 -> 0, -1 -> 1 (correct) & match \\
Codex, Trial 3 & +1 -> 0, -1 -> 1 (correct) & match \\
\bottomrule
\end{tabular}
\end{center}

\textbf{Code Snippets:}

\emph{Claude Code, Trial 3 --- incorrect mapping; comment documents the wrong assumption:}

\begin{minted}[bgcolor=codebg, fontsize=\footnotesize, breaklines, frame=none, xleftmargin=0.5em]{python}
# Choice: -1 (left) -> 0, 1 (right) -> 1
choice = ((choice_raw + 1) / 2).astype(np.int32)   # {-1, 1} -> {0, 1}
...
output_trial[0, :] = choice[trial_idx]             # broadcast across 100 bins
\end{minted}

The comment captures the agent's mental model: "-1 = left". The arithmetic \texttt{(x + 1) / 2} then collapses \texttt{\{-1, +1\}} to \texttt{\{0, 1\}} in that order, so every per-trial label ends up on the wrong side. The other two Claude Code trials reach the same flipped result by slightly different code (e.g. \texttt{np.where(choice == 1, 1, 0)}).

\emph{Codex, Trial 2 --- explicit per-value mapping that matches the IBL convention:}

\begin{minted}[bgcolor=codebg, fontsize=\footnotesize, breaklines, frame=none, xleftmargin=0.5em]{python}
def map_choice(raw_choice):
    mapped = np.empty(raw_choice.shape[0], dtype=np.int8)
    mapped[raw_choice ==  1] = 0   # left
    mapped[raw_choice == -1] = 1   # right
    return mapped
...
out = np.vstack([
    np.full(NBINS, choice[trial_idx], dtype=np.int8),
    ...
])
\end{minted}

Each raw value is mapped explicitly with an inline left/right comment that matches the \texttt{ibllib} convention. Reading the IBL library code rather than guessing from the sign of the integer is what separates this trial from the three Claude Code runs.

\end{agentexample}

\begin{agentexample}[Allen2P, \texttt{ASSUME/VARNAME}]
\textbf{Task:} The agent must group the experiment table into output sessions. The Allen Brain Observatory Visual Behavior 2P dataset uses two acquisition rigs --- a Scientifica single-plane scope (1 imaging plane per session) and a Multiscope multi-plane scope (up to 8 simultaneously imaged planes per session). The dataset's own \texttt{methods.txt} is unambiguous: \emph{"The data collected in a single continuous recording is defined as a session. For multi-plane imaging experiments, there can be up to 8 imaging planes (8 experiments) per session."} Each plane is packaged as a separate NWB file \texttt{behavior\_ophys\_experiment\_<id>.nwb}, and the experiment table carries both \texttt{ophys\_experiment\_id} (one row per plane) and \texttt{ophys\_session\_id} (shared across planes from the same simultaneous recording). The \textbf{nominally correct} decision is to group experiments by \texttt{ophys\_session\_id} --- e.g. \texttt{mouse\_exps.drop\_duplicates(subset='ophys\_session\_id')} --- so that one Multiscope session ends up as a single output entry containing all of its planes. Treating each \texttt{ophys\_experiment\_id} as its own session splits one behavioural recording across up to 8 separate output entries and breaks any analysis that depends on across-plane neuron co-occurrence within a session.

\textbf{Summary:} All six trials produced the same wrong answer: each iterated the experiment table row-by-row and emitted one output session per \texttt{ophys\_experiment\_id}. The striking part is that every agent's own \texttt{CONVERSION\_NOTES.md} explicitly \emph{acknowledged} the experiment / session distinction before then choosing to ignore it. Codex Trial 3's note is the loudest: \emph{"the conversion intentionally treats each \texttt{ophys\_experiment\_id} plane as a separate session."} This is therefore both a VARNAME failure (in most datasets \texttt{experiment\_id} and \texttt{session\_id} are interchangeable, so the names invite the wrong shortcut) and an ASSUME failure (the agent saw the dataset's note and the AllenSDK tutorial code yet defaulted to the common-case mapping anyway).

\begin{center}
\begin{tabular}{lll}
\toprule
Trial & Per-experiment key line & Rating \\
\midrule
\rowcolor{red!20} Claude Code, Trial 1 & \texttt{exp\_id = row['ophys\_experiment\_id']} & incorrect \\
\rowcolor{red!20} Claude Code, Trial 2 & \texttt{eid = row['ophys\_experiment\_id']} & incorrect \\
\rowcolor{red!20} Claude Code, Trial 3 & \texttt{eid = row['ophys\_experiment\_id']} & incorrect \\
\rowcolor{red!20} Codex, Trial 1 & \texttt{SessionPreview(experiment\_id=experiment\_id, ...)} & incorrect \\
\rowcolor{red!20} Codex, Trial 2 & \texttt{SessionInfo(ophys\_experiment\_id=row.ophys\_experiment\_id, ...)} & incorrect \\
\rowcolor{red!20} Codex, Trial 3 & \texttt{SessionMeta(ophys\_experiment\_id=row.ophys\_experiment\_id, ...)} & incorrect \\
\bottomrule
\end{tabular}
\end{center}

\textbf{Code Snippets:}

\emph{Claude Code, Trial 2 --- per-\texttt{ophys\_experiment\_id} loop, no session grouping:}

\begin{minted}[bgcolor=codebg, fontsize=\footnotesize, breaklines, frame=none, xleftmargin=0.5em]{python}
active_exps = our_exps[our_exps['passive'] == False].copy()
active_exps = active_exps.sort_values('ophys_experiment_id').reset_index(drop=True)

for idx, (_, row) in enumerate(active_exps.iterrows()):
    eid       = row['ophys_experiment_id']
    nwb_path  = nwb_map[eid]
    result    = process_experiment(nwb_path, row, collect_stats_only=False)
\end{minted}

The experiment table is filtered for active sessions and sorted by \texttt{ophys\_experiment\_id}; the loop body then processes one NWB file at a time and emits one output entry per file. Multi-plane sessions are silently fragmented into their constituent planes. The other two Claude Code trials follow the same pattern with minor variations.

\emph{Codex, Trial 3 --- same per-experiment loop, with the deliberate-but-wrong design choice in the agent's notes:}

\begin{minted}[bgcolor=codebg, fontsize=\footnotesize, breaklines, frame=none, xleftmargin=0.5em]{python}
exp_table = exp_table.sort_values("ophys_experiment_id")

sessions = []
for row in exp_table.itertuples(index=False):
    sessions.append(SessionMeta(
        ophys_experiment_id  = int(row.ophys_experiment_id),
        ophys_session_id     = int(row.ophys_session_id),     # read but unused
        behavior_session_id  = int(row.behavior_session_id),  # read but unused
        mouse_id             = str(row.mouse_id),
        ...
    ))
\end{minted}

The agent reads both \texttt{ophys\_experiment\_id} \emph{and} \texttt{ophys\_session\_id} into its \texttt{SessionMeta}, then keys output sessions on \texttt{ophys\_experiment\_id} regardless. Its accompanying \texttt{CONVERSION\_NOTES.md} adds: \emph{"the conversion intentionally treats each \texttt{ophys\_experiment\_id} plane as a separate session"} --- the wrong granularity was a self-aware design choice, not an oversight.

\emph{Reference solution --- group experiments by \texttt{ophys\_session\_id}, then enumerate per-session planes:}

\begin{minted}[bgcolor=codebg, fontsize=\footnotesize, breaklines, frame=none, xleftmargin=0.5em]{python}
mouse_exps     = vb_experiments[vb_experiments.mouse_id == mouse_id]
mouse_sessions = mouse_exps.drop_duplicates(subset='ophys_session_id')[
    ['ophys_session_id', 'session_type', 'date_of_acquisition']
].sort_values(by='date_of_acquisition')

for sid in mouse_sessions['ophys_session_id'].values:
    sess_exps = mouse_exps[mouse_exps.ophys_session_id == sid]   # 1 to 8 planes
    # ... build one output session from sess_exps (combine planes' neurons,
    #     reuse the shared behavioural data) ...
\end{minted}

The canonical grouping: one row per \texttt{ophys\_session\_id}, ordered chronologically. For each session, \texttt{sess\_exps} enumerates the simultaneously imaged planes --- one row for Scientifica, up to eight for Multiscope. A single output entry is then assembled from those planes (neuron rows concatenated across planes, behavioural data taken once since it is shared).

\end{agentexample}
